%%% ---------------------------------------------------------------------------
%%% preamble.tex —— 你自己的宏包与自定义命令
%%%
%%% 类文件里已经载入了：
%%%   ctex / fontspec / geometry / unicode-math / acadmath / amsmath / mathtools
%%%   graphicx / booktabs / array / multirow / makecell / longtable / threeparttable
%%%   caption / subcaption / float / algorithm / algpseudocode / amsthm
%%%   biblatex / fancyhdr / xcolor / hyperref / bookmark / cleveref
%%%   enumitem / url / microtype
%%% 不要重复载入这些，尤其不要重新载入 hyperref。
%%%
%%% 共享的数学记号（\vect \mat \norm \argmin \dd \subtext …）见 ../common/acadmath.sty，
%%% 那个文件同时被 beamer 模板使用，改它可以让论文和报告的符号保持一致。
%%% ---------------------------------------------------------------------------

%%% ---- 按需追加的宏包 ----
% \usepackage{tikz}
% \usepackage{pgfplots}\pgfplotsset{compat=1.18}
% \usepackage{listings}          % 代码
% \usepackage{physics}           % 注意：与 unicode-math 有已知冲突，慎用

%%% ---- 本文专用的记号 ----
% 例：把常用的模型名、数据集名定成命令，避免全文写法不一致
\newcommand{\ourmethod}{MSFF-Net}
\newcommand{\dataset}{DOTA-v2.0}

% 例：本文专用的数学量
\newcommand{\featmap}{\mat{F}}          % 特征图
\newcommand{\loss}{\mathcal{L}}         % 损失函数
\newcommand{\params}{\vect{\theta}}     % 参数向量

%%% ---- 占位图 ----
%%% 只是为了让模板在没有任何图片文件时也能编译。
%%% 真正插图时删掉它，改用：\includegraphics[width=0.85\linewidth]{architecture}
%%% 用法：\placeholder{<宽>}{<高>}{<说明文字>}
\newcommand{\placeholder}[3]{%
  \begingroup
    \fboxsep=0pt \fboxrule=0.4pt
    \fbox{\parbox[c][#2][c]{\dimexpr#1-2\fboxrule\relax}{\centering\footnotesize\ttfamily\color{gray}#3}}%
  \endgroup
}

%%% ---- 若要覆盖类文件里的样式，写在这里（会后于类文件生效）----
% 例：改正文字号
% \AtBeginDocument{\zihao{-4}}
%
% 例：图题改成左对齐
% \captionsetup[figure]{justification=raggedright, singlelinecheck=false}
